% =============================================================
%  Kapitel 1 – Einleitung
% =============================================================
\section{Einleitung}
\label{sec:einleitung}

Dieses Dokument beschreibt das \textbf{Program Design} des verteilten Not-Halt-Systems
einer autonomen Fahrzeugplattform (Truck).
Es hat den Charakter eines \ac{SDD} und dokumentiert
die Softwarearchitektur sowie das detaillierte Design des Teilsystems \textbf{Emergency Stop}.

\subsection{Scope}
\label{sec:scope}

\begin{itemize}
  \item \textbf{In Scope:} Software im Truck (Jetson Nano, Arduino Mega Master,
        Arduino Mega Slave), Architektur, Kommunikationsschnittstellen,
        Emergency-Stop-Logik, Boundary Conditions.
  \item \textbf{Out of Scope:} PC-Anwendung, \ac{ROS}-Integration, \ac{WLAN}-Kommunikation,
        Details der Sensor-Fusion (wird als Blackbox behandelt).
  \item \textbf{Vereinfachungen:}
        LiDAR liefert direkt die Distanz zum nächstliegenden Objekt;
        Sensor-Fusion ist vorgelagert und für dieses Dokument transparent.
\end{itemize}

\subsection{Ziel des Dokuments}
\label{sec:ziel}

Das \ac{SDD} richtet sich an Entwicklerinnen und Entwickler, die
\begin{enumerate}
  \item die Gesamtarchitektur und die Subsystem-Verantwortlichkeiten verstehen,
  \item das Deployment auf den Hardwareeinheiten nachvollziehen,
  \item das dynamische Verhalten (Sequenzen, Boundary Conditions) analysieren und
  \item das Emergency-Stop-Modul detailliert implementieren wollen.
\end{enumerate}

\newpage
\subsection{Abgrenzung Normal Operation vs.\ Emergency Stop}
\label{sec:abgrenzung}

Das System unterscheidet zwei Betriebsmodi (vgl.\ Abbildung~\ref{fig:zone_definition}):


\begin{description}
  \item[Normal Operation (Gelb/Grün)]
    Der \emph{\ac{TTC}-Observer} überwacht die Umgebung und erlaubt
    eine automatische Reaktion auf Hindernisse.  Der Fahrvorgang wird angepasst,
    das Fahrzeug kann aber weiterfahren.  \ac{TTC} ist \textbf{kein} Bestandteil des
    Emergency-Stop-Moduls.
  \item[Safety Zone / Emergency Stop (Rot)]
    Sobald ein Objekt die rote Sicherheitszone unterschreitet, ein Pilzknopf gedrückt
    oder ein Fronttaster ausgelöst wird, übernimmt der \emph{Emergency Stop}.
    Das Fahrzeug bleibt stehen, bis eine \textbf{manuelle Freigabe} erfolgt
    (Guard-Bedingung: Safety-Zone nicht mehr aktiv, kein Fehler).
\end{description}


\begin{figure}[H]
  \centering
  \includegraphics[width=0.95\textwidth]{Zone_Definition.png}
  \caption{Zone Definition (Ref. Aufgabenstellung \cite{PPT_SE_Operation_Zones})}
  \label{fig:zone_definition}
\end{figure}

